\newpage

\namedgridquestion{14}{
Escriba una función en Python que calcule el \textbf{saldo anterior} de un cliente.
La función debe recibir como parámetros la cuota mensual, el gasto actual y el número de pagos.
Debe retornar el saldo anterior.

La cuota mensual se calcula como:

\[
\text{cuota} =
\text{gasto\_actual}
+
\frac{\text{saldo\_anterior} \cdot i \cdot (1+i)^n}
{(1+i)^n - 1}
\]

Donde:

\[
i = 0.025 \qquad \text{y} \qquad n = \text{numero\_pagos}
\]

Considere que esta función pertenece a la capa de lógica; ergo, no interactúa con el usuario.
}
